% GENERATED FILE. Edit canonical lessons, then run npm run generate:books.
% canonical-source-hash: fnv1a64:28657bdbbd039572
% canonical-lessons: GU-C71-white, GU-C71-brown, GU-C71-pink, GU-C71-beautiful, GU-C71-bad

\chapter{White, Brown, Pink, and What Looks Good}
\label{ch:gu-white-brown-pink-and-what-looks-good}
\hlchaptermodality{71}

\begin{chapteropening}
\textbf{By the end of this chapter:} \emph{I can name white, brown and pink, and say beautiful or bad.}

Complete the last lesson of chapter 71: I can name white, brown and pink, and say beautiful or bad.
\end{chapteropening}

\section[\texorpdfstring{\gu{સફેદ} (saphed)}{સફેદ (saphed)}]{\gu{સફેદ} (saphed) — white}

\label{lesson:GU-C71-white}



\begin{quote}
Before the new one: say the Gujarati for blue, then the Gujarati for black.
\end{quote}

\subsection*{You'll want to know: \gu{સફેદ}}
\textbf{\gu{સફેદ}} (\emph{saphed}) — \textquotedblleft{}white\textquotedblright{}.

\begin{cousinweb}
\textbf{\gu{સફેદ}} (\emph{saphed}) is a loan, usually traced to Persian \emph{safēd}, \textquotedblleft{}white\textquotedblright{}.
\end{cousinweb}

\begin{grammarlens}[title={What you've built}]
One more word for white, brown, pink, and what looks good.
\end{grammarlens}

\subsection*{Your turn}
\begin{itemize}
\raggedright
  \item \emph{Say it:} \emph{saphed}
  \item \emph{Say it:} \emph{saphed}, once more
  \item \emph{Say it:} say \emph{saphed}
  \item \emph{Recall:} say the Gujarati for blue, then the Gujarati for black, then say \emph{saphed} again
\end{itemize}

\begin{tcolorbox}[breakable,colback=teal!4,colframe=teal!35!black,title={Before you move on}]
What does \gu{સફેદ} mean? (\textquotedblleft{}White\textquotedblright{}.) Say it once more.
\end{tcolorbox}
\section[\texorpdfstring{\gu{ભૂરો} (bhūro)}{ભૂરો (bhūro)}]{\gu{ભૂરો} (bhūro) — brown}

\label{lesson:GU-C71-brown}



\begin{quote}
Before the new one: say the Gujarati for black, then the Gujarati for white.
\end{quote}

\subsection*{You'll want to know: \gu{ભૂરો}}
\textbf{\gu{ભૂરો}} (\emph{bhūro}) — \textquotedblleft{}brown\textquotedblright{}.

\begin{cousinweb}
\textbf{\gu{ભૂરો}} (\emph{bhūro}) is a tadbhava. It comes from \emph{*bhrūra}, \textquotedblleft{}brown\textquotedblright{}, a form CDIAL reconstructs and sets beside Indo-European \emph{*bhrūro-} and \emph{*bhrūno-} (CDIAL 9690). The second is the source of English \emph{brown}.
\end{cousinweb}

\begin{grammarlens}[title={What you've built}]
One more word for white, brown, pink, and what looks good.
\end{grammarlens}

\subsection*{Your turn}
\begin{itemize}
\raggedright
  \item \emph{Say it:} \emph{bhūro}
  \item \emph{Say it:} \emph{bhūro}, once more
  \item \emph{Say it:} say \emph{bhūro}
  \item \emph{Recall:} say the Gujarati for black, then the Gujarati for white, then say \emph{bhūro} again
\end{itemize}

\begin{tcolorbox}[breakable,colback=teal!4,colframe=teal!35!black,title={Before you move on}]
What does \gu{ભૂરો} mean? (\textquotedblleft{}Brown\textquotedblright{}.) Say it once more.
\end{tcolorbox}
\section[\texorpdfstring{\gu{ગુલાબી} (gulābī)}{ગુલાબી (gulābī)}]{\gu{ગુલાબી} (gulābī) — pink}

\label{lesson:GU-C71-pink}



\begin{quote}
Before the new one: say the Gujarati for white, then the Gujarati for brown.
\end{quote}

\subsection*{You'll want to know: \gu{ગુલાબી}}
\textbf{\gu{ગુલાબી}} (\emph{gulābī}) — \textquotedblleft{}pink\textquotedblright{}.

\begin{cousinweb}
\textbf{\gu{ગુલાબી}} (\emph{gulābī}) is a loan, usually traced to Persian \emph{gulābī}, \textquotedblleft{}rose-coloured\textquotedblright{}, from \emph{gulāb}, \textquotedblleft{}rose-water\textquotedblright{}.
\end{cousinweb}

\begin{grammarlens}[title={What you've built}]
One more word for white, brown, pink, and what looks good.
\end{grammarlens}

\subsection*{Your turn}
\begin{itemize}
\raggedright
  \item \emph{Say it:} \emph{gulābī}
  \item \emph{Say it:} \emph{gulābī}, once more
  \item \emph{Say it:} say \emph{gulābī}
  \item \emph{Recall:} say the Gujarati for white, then the Gujarati for brown, then say \emph{gulābī} again
\end{itemize}

\begin{tcolorbox}[breakable,colback=teal!4,colframe=teal!35!black,title={Before you move on}]
What does \gu{ગુલાબી} mean? (\textquotedblleft{}Pink\textquotedblright{}.) Say it once more.
\end{tcolorbox}
\section[\texorpdfstring{\gu{સુંદર} (sundar)}{સુંદર (sundar)}]{\gu{સુંદર} (sundar) — beautiful}

\label{lesson:GU-C71-beautiful}



\begin{quote}
Before the new one: say the Gujarati for brown, then the Gujarati for pink.
\end{quote}

\subsection*{You'll want to know: \gu{સુંદર}}
\textbf{\gu{સુંદર}} (\emph{sundar}) — \textquotedblleft{}beautiful\textquotedblright{}.

\begin{cousinweb}
\textbf{\gu{સુંદર}} (\emph{sundar}) is Sanskrit \emph{sundara}, \textquotedblleft{}beautiful\textquotedblright{}, taken over whole as a learned loan, a tatsama (CDIAL 13474). CDIAL asks whether the Sanskrit word is itself a Middle Indo-Aryan form of \emph{sūnara}, \textquotedblleft{}excellent\textquotedblright{}.
\end{cousinweb}

\begin{grammarlens}[title={What you've built}]
One more word for white, brown, pink, and what looks good.
\end{grammarlens}

\subsection*{Your turn}
\begin{itemize}
\raggedright
  \item \emph{Say it:} \emph{sundar}
  \item \emph{Say it:} \emph{sundar}, once more
  \item \emph{Say it:} say \emph{sundar}
  \item \emph{Recall:} say the Gujarati for brown, then the Gujarati for pink, then say \emph{sundar} again
\end{itemize}

\begin{tcolorbox}[breakable,colback=teal!4,colframe=teal!35!black,title={Before you move on}]
What does \gu{સુંદર} mean? (\textquotedblleft{}Beautiful\textquotedblright{}.) Say it once more.
\end{tcolorbox}
\section[\texorpdfstring{\gu{ખરાબ} (kharāb)}{ખરાબ (kharāb)}]{\gu{ખરાબ} (kharāb) — bad}

\label{lesson:GU-C71-bad}



\begin{quote}
Before the new one: say the Gujarati for pink, then the Gujarati for beautiful.
\end{quote}

\subsection*{You'll want to know: \gu{ખરાબ}}
\textbf{\gu{ખરાબ}} (\emph{kharāb}) — \textquotedblleft{}bad\textquotedblright{}. The opposite of \textbf{\gu{સારું}}.

\begin{cousinweb}
\textbf{\gu{ખરાબ}} (\emph{kharāb}) is a loan from Arabic \emph{kharāb}, \textquotedblleft{}ruin, desolation\textquotedblright{}, which reached Gujarati through Persian. Its Arabic root is \emph{kh-r-b} (Lane).
\end{cousinweb}

\begin{grammarlens}[title={What you've built}]
One more word for white, brown, pink, and what looks good.
\end{grammarlens}

\subsection*{Your turn}
\begin{itemize}
\raggedright
  \item \emph{Say it:} \emph{kharāb}
  \item \emph{Say it:} \emph{kharāb}, once more
  \item \emph{Say it:} say \emph{kharāb}
  \item \emph{Recall:} say the Gujarati for pink, then the Gujarati for beautiful, then say \emph{kharāb} again
\end{itemize}

\begin{tcolorbox}[breakable,colback=teal!4,colframe=teal!35!black,title={Before you move on}]
What does \gu{ખરાબ} mean? (\textquotedblleft{}Bad\textquotedblright{}.) Say it once more.
\end{tcolorbox}
